% !TEX root = ../../../main.tex
\section{Diagonalization}
\label{sec:linear-algebra-i:diagonalization}

% Sources: IMG_9268, IMG_9323--IMG_9324, IMG_9328;
% October 15 PDF, pp. 22--26.
\begin{definition}
  A square matrix \(M\) is diagonalizable over \(\mathbb F\) if it is
  similar to a diagonal matrix. Equivalently, there exist an invertible
  matrix \(S\) and a diagonal matrix \(D\) such that
  \[
    M=SDS^{-1}.
  \]
\end{definition}

\begin{theorem}[Diagonalization and eigenbases]
  \label{thm:linear-algebra-i:diagonalization-eigenbasis}
  A matrix \(M\in\mathbb F^{n\times n}\) is diagonalizable if and only
  if \(\mathbb F^n\) has a basis of eigenvectors of \(M\).
\end{theorem}

\begin{marginfigure}
  \centering
  % MS=SD: S converts eigenbasis coordinates into the original coordinates.
\begin{tikzpicture}[
  x=1mm,y=1mm,font=\footnotesize,line width=0.45pt,
  every node/.style={inner sep=3pt},
  map/.style={-{Latex[length=1.6mm,width=1.1mm]}}
]
  \node (c) at (0,24) {\(c\)};
  \node (Dc) at (32,24) {\(Dc\)};
  \node (Sc) at (0,0) {\(Sc\)};
  \node (SDc) at (32,0) {\(SDc\)};
  \draw[map] (c) -- node[above] {\(D\)} (Dc);
  \draw[map] (c) -- node[left] {\(S\)} (Sc);
  \draw[map] (Dc) -- node[right] {\(S\)} (SDc);
  \draw[map] (Sc) -- node[below] {\(M\)} (SDc);
  \path (current bounding box.south) ++(0,-2mm);
\end{tikzpicture}

  \caption[Diagonalization through coordinates]{The top row acts in
    eigenbasis coordinates; the bottom acts in the original coordinates.
    Both routes give \(MSc=SDc\).}
  \label{fig:linear-algebra-i:diagonalization-coordinates}
\end{marginfigure}

\begin{proof}
  Suppose \(M=SDS^{-1}\), where \(D\) has diagonal entries
  \(\lambda_1,\ldots,\lambda_n\). Let \(c_k=Se_k\) be column \(k\)
  of \(S\). Since \(S\) is invertible, its columns form a basis, and
  \begin{align*}
    Mc_k&=MSe_k=SDS^{-1}Se_k\\
      &=SDe_k=S(\lambda_ke_k)=\lambda_kc_k.
  \end{align*}
  Thus each column is an eigenvector.

  Conversely, suppose \((b_1,\ldots,b_n)\) is a basis of eigenvectors,
  with \(Mb_k=\lambda_kb_k\). Let \(S\) have these basis vectors as
  columns, and let \(D\) have the corresponding eigenvalues on its diagonal.
  The column equations say \(MS=SD\). Since \(S\) is invertible,
  \(M=SDS^{-1}\).
\end{proof}

The identity \(MS=SD\) is visualized in
\autoref{fig:linear-algebra-i:diagonalization-coordinates}.
If \(c\) is an eigenbasis coordinate column, \(S\) reconstructs the vector
\(Sc\), while \(D\) multiplies each coordinate by its eigenvalue.

\begin{definition}
  An endomorphism \(L:V\to V\) is diagonalizable if \(V\) has a basis
  \((e_1,\ldots,e_n)\) of eigenvectors of \(L\).
\end{definition}

For such a basis,
\[
  L(e_k)=\lambda_ke_k,\qquad
  L(E)=E\begin{bmatrix}
    \lambda_1&0&\cdots&0\\
    0&\lambda_2&\cdots&0\\
    \vdots&\vdots&\ddots&\vdots\\
    0&0&\cdots&\lambda_n
  \end{bmatrix}.
\]
Equivalently, the matrix representing \(L\) in any basis is similar
to a diagonal matrix.

\subsection{Failure of diagonalization}

\begin{example}
  Let \((e_1,e_2)\) be a basis of \(V\), and define
  \[
    L(e_1)=e_2,\qquad L(e_2)=0.
  \]
  For \(v=ae_1+be_2\), \(L(v)=ae_2\). The eigenvector equation is
  \[
    ae_2=\lambda(ae_1+be_2).
  \]
  If \(\lambda\ne0\), comparison of coefficients gives \(a=0\)
  and then \(b=0\), so there is no nonzero eigenvector. For
  \(\lambda=0\), it gives \(a=0\). Thus the only eigenvalue is zero
  and its eigenspace is \(\operatorname{span}\{e_2\}\). It cannot
  provide a basis of \(V\), so \(L\) is not diagonalizable.
\end{example}

The nonzero matrix \(N\) in
\autoref{ex:linear-algebra-i:same-polynomial-not-similar} similarly has
only the eigenvalue zero and a one-dimensional eigenspace, so it is not
diagonalizable. The rotation matrix in
\autoref{ex:linear-algebra-i:rotation-no-real-eigenvalue} is not
diagonalizable over \(\R\), since it has no real eigenvectors.

% !TEX root = ../../hw06.tex
% Source: supplied Homework 6 PDF, page 2, Question 6.
% Retain the source's "zero matrix" wording for the linear transformation;
% its unstated nonzero-space hypothesis is recorded in the README.
\begin{exercise}
  \label{ex:hw06:nilpotent-eigenvalues}
  A linear transformation \(A:V\to V\) is called \emph{nilpotent} if for some
  positive integer \(k\), \(A^k\) is the zero matrix. Prove that the only
  eigenvalue of \(A\) is \(0\).
  \solutionlink{hw06:nilpotent-eigenvalues}
\end{exercise}

